\chapter{幂}\label{ch:g8:powers}

将一张纸对折 $10$ 次（如果你能做到！），它的厚度
每次都会加倍，$10$ 次加倍使其厚度变为原来的
$2^{10} = 1024$ 倍。幂是重复乘法的简写；
本章介绍幂的记号及其运算法则，
并特别强调 $10$ 的幂的特殊作用。

\section{定义}

\begin{definition}[幂]\label{def:g8:powers:def}
设数 $a$ 与整数 $n \geq 1$：
\[
a^n = \underbrace{a \times a \times \dots \times a}_{n
\text{ 个因数}},
\]
读作“$a$ 的 $n$ 次幂”；$a$ 为\emph{底数}，$n$ 为
\emph{指数}\index{指数}。特别地，$a^2$ 读作“$a$
的平方”，$a^3$ 读作“$a$ 的立方”。约定如下：
\[
a^1 = a,
\qquad
a^0 = 1 \ (a \neq 0),
\qquad
a^{-n} = \frac{1}{a^n}.
\]
\end{definition}

\begin{example}\label{ex:g8:powers:def}
$3^4 = 3 \times 3 \times 3 \times 3 = 81$; \quad
$10^3 = 1000$; \quad $5^{-2} = \frac{1}{25} = 0.04$; \quad
$(-2)^3 = -8$ 与 $(-2)^4 = +16$
（见\cref{ex:g8:negprod:powers}）。注意：$a^n$ 并\emph{不}等于
$a \times n$：\ $2^5 = 32$，而不是 $10$。
\end{example}

\begin{omfigure}
\begin{tikzpicture}[scale=0.9]
  \foreach \n/\x in {0/0, 1/1.4, 2/2.8, 3/4.2, 4/5.6} {
    \pgfmathsetmacro{\h}{0.28*pow(2,\n)}
    \draw[omDef, fill=omDef!20] (\x,0) rectangle (\x+0.9,\h);
    \node[below, font=\small] at (\x+0.45,0) {$2^\n$};
    \pgfmathtruncatemacro{\v}{pow(2,\n)}
    \node[above, font=\small] at (\x+0.45,\h) {$\v$};
  }
\end{tikzpicture}

{\small $2$ 的幂：每根柱子都是前一根的两倍。
通过重复相乘实现的增长
远快于通过重复相加实现的增长。}
\end{omfigure}

\section{幂的运算法则}

\begin{theorem}[幂的运算法则]\label{thm:g8:powers:rules}
设非零底数 $a$，整数\omterm{def:g8:powers:def}{指数} $m$、$n$：
\[
a^m \times a^n = a^{m+n},
\qquad
\frac{a^m}{a^n} = a^{m-n},
\qquad
\left(a^m\right)^n = a^{m \times n},
\qquad
(ab)^n = a^n b^n .
\]
\end{theorem}

\begin{proof}[通过计算因数个数证明]
$a^m \times a^n$ 是把 $m$ 个 $a$ 与其后的 $n$ 个\omterm{def:g5:division:multiple}{因数}排成一列：
共 $m + n$ 个\omterm{def:g5:division:multiple}{因数}。$\left(a^m\right)^n$ 是把由 $m$ 个\omterm{def:g5:division:multiple}{因数}组成的一段重复
$n$ 次，共 $mn$ 个\omterm{def:g5:division:multiple}{因数}。$(ab)^n$ 含 $n$ 个字母 $a$ 与 $n$
个字母 $b$，可以重新组合。商的法则来自于约去
$n$ 个\omterm{def:g5:division:multiple}{因数}（共 $m$ 个）；若约定
$a^0 = 1$ 与 $a^{-n} = \frac{1}{a^n}$，它
在 $n \geq m$ 时仍然成立。
\end{proof}

\begin{example}\label{ex:g8:powers:rules}
\[
2^3 \times 2^5 = 2^8 = 256,
\qquad
\frac{7^6}{7^4} = 7^2 = 49,
\qquad
\left(5^2\right)^3 = 5^6,
\qquad
2^4 \times 5^4 = 10^4 .
\]
注意：运算法则只适用于\emph{同底数}的情形（或
同\omterm{def:g8:powers:def}{指数}的情形，对最后一个式子而言）。没有任何法则可以化简 $2^3 \times 5^2$——
只需直接计算：$8 \times 25 = 200$。
\end{example}

\section{10 的幂}

\begin{proposition}[10 的幂]\label{prop:g8:powers:ten}
对 $n \geq 1$，有 $10^n = 1\underbrace{0\dots0}_{n}$，且
$10^{-n} = 0.\underbrace{0\dots0}_{n-1}1$。幂的运算法则
表明：10 的幂相乘时，\omterm{def:g8:powers:def}{指数}相加。
\end{proposition}

\begin{example}\label{ex:g8:powers:ten}
$10^4 \times 10^3 = 10^7$; \quad
$\dfrac{10^2}{10^5} = 10^{-3} = 0.001$。大量与微小量
都变得易于读写：
\[
\text{十亿} = 10^9,
\qquad
\text{百万分之一} = 10^{-6}.
\]
结合\omterm{def:g6:decimals:places}{小数}：$3.2 \times 10^5 = 320\,000$ 与
$4.7 \times 10^{-3} = 0.0047$（移动\omterm{def:g4:decimals:point}{小数点}，
见\cref{prop:g6:decimals:shift}）。这种写法的系统应用——
\emph{科学记数法}——将
在\cref{ch:g9:fractions}中展开。
\end{example}

\begin{example}[数量级]\label{ex:g8:powers:magnitude}
光速约为 $3 \times 10^8$ m/s；一年约有
$3.2 \times 10^7$ 秒。因此一光年约为
\[
3 \times 3.2 \times 10^{8+7} \approx 10 \times 10^{15} = 10^{16}
\text{ m}
\]
——即一亿亿米。借助十的幂，天文计算
可以写在一行之内。
\end{example}

\begin{method}[化简含幂的表达式]\label{met:g8:powers:simplify}
\begin{enumerate}
  \item 按底数归并各因式；
  \item 对每个底数应用\omterm{def:g8:powers:def}{指数}运算法则；
  \item 计算其余较小的幂，若结果较大，则保留
        为幂的形式。
\end{enumerate}
\end{method}

\begin{example}\label{ex:g8:powers:simplify}
\[
\frac{3^5 \times 3^{-2} \times 4^2}{3^2}
= \frac{3^{5 + (-2)}}{3^2} \times 16
= 3^{3 - 2} \times 16
= 3 \times 16 = 48 .
\]
\end{example}

\section{练习}

\begin{exercise}[$\star$]\label{exo:g8:powers:1}
计算：
\[
2^6, \qquad 3^3, \qquad 10^5, \qquad 1^{100}, \qquad 0.1^2, \qquad
6^0 .
\]
\end{exercise}

\begin{exercise}[$\star$]\label{exo:g8:powers:2}
计算：
\[
(-3)^2, \qquad -3^2, \qquad (-1)^{15}, \qquad (-5)^3, \qquad
\left(\tfrac{2}{3}\right)^2 .
\]
\end{exercise}

\begin{exercise}[$\star$]\label{exo:g8:powers:3}
将下列各式写成一个幂：
\[
7^4 \times 7^5, \qquad
\frac{2^9}{2^3}, \qquad
\left(10^3\right)^4, \qquad
5^6 \times 5^{-2}, \qquad
3^4 \times 7^4 .
\]
\end{exercise}

\begin{exercise}[$\star$]\label{exo:g8:powers:4}
写成\omterm{def:g6:decimals:places}{小数}：$10^{-2}$；\quad $4 \times 10^3$；\quad
$2.5 \times 10^{-4}$; \quad $10^0$.
\end{exercise}

\begin{exercise}[$\star$]\label{exo:g8:powers:5}
写成十的幂的形式：十万；十分之一；一
百亿；$0.000\,001$。
\end{exercise}

\begin{exercise}[$\star$]\label{exo:g8:powers:6}
化简后计算：
\[
\frac{10^7 \times 10^{-3}}{10^2},
\qquad
\frac{2^5 \times 2^4}{2^6},
\qquad
\left(2^2\right)^3 \times 2^{-4} .
\]
\end{exercise}

\begin{exercise}[$\star$]\label{exo:g8:powers:7}
判断下列各式是否正确？并改正错误的式子。
\[
2^3 \times 2^4 = 2^{12}; \qquad
5^2 + 5^3 = 5^5; \qquad
(3^2)^4 = 3^8; \qquad
10^3 \times 10^3 = 100^3 .
\]
\end{exercise}

\begin{exercise}[$\star\star$]\label{exo:g8:powers:8}
一则谣言传播：第 1 天有 3 人知道；每天，每个
知道的人每人再告诉 3 个新的人。用幂表示
第 $4$ 天被告知的\emph{新}人数，并求
出第 4 天结束时知道谣言的人数（包括最初的 3 人）。
\end{exercise}

\begin{exercise}[$\star\star$]\label{exo:g8:powers:9}
一张纸厚 $0.1$ mm，即 $10^{-4}$ m。每折一次，
厚度便加倍。
\begin{enumerate}
  \item 将折叠 $10$ 次后的厚度写成乘积形式，并
        以厘米为单位计算该厚度（$2^{10} = 1024$）。
  \item 折叠 $42$ 次后，厚度为
        $2^{42} \times 10^{-4}$ m，其中 $2^{42} \approx 4.4 \times
        10^{12}$。证明该厚度超过地月距离，
        约为 $3.8 \times 10^8$ m。
\end{enumerate}
\end{exercise}

\begin{exercise}[$\star\star$]\label{exo:g8:powers:10}
不用计算器，按从小到大的顺序排列：
\[
2^{10}, \qquad 10^3, \qquad 3^6, \qquad 5^4 .
\]
（逐一计算；$2^{10}$ 与 $10^3$ 是著名的“近邻”。）
\end{exercise}

\begin{exercise}[$\star\star\star$]\label{exo:g8:powers:11}
哪个更大，$2^{100}$ 还是 $10^{30}$？利用
$2^{10} = 1024 > 10^3$，比较
$2^{100} = \left(2^{10}\right)^{10}$ 与
$\left(10^3\right)^{10}$。
\end{exercise}

\section{问题：棋盘与二的幂}

\begin{problem}[{周末问题——等比数列求和
$1 + 2 + 4 + \dots + 2^{n-1} = 2^n - 1$，由一个著名传说通向
二进制数}]
\label{pb:g8:powers:1}
传说：作为发明国际象棋的奖赏，贤者西萨
请求国王在棋盘第一格放一粒麦粒，
第二格放两粒，第三格放四粒——从一格到下一格都翻倍，
直到第六十四格。国王却为这样的“谦逊”大笑。
本题计算国王许下的承诺，
利用\cref{thm:g8:powers:rules}中的\omterm{def:g8:powers:def}{指数}法则，并
最终落到这个故事暗中指向的地方：每台计算机内部的二进制数。

\medskip
\noindent\textbf{第一部分——翻倍技巧。}
对 $n \geq 1$，设 $S_n$ 为前
$n$ 格上的麦粒总数。
\begin{enumerate}
  \item 将第 $k$ 格上的麦粒数写成
        $2$ 的幂。第 $64$ 格上对应哪一个幂？
  \item 计算 $S_1$、$S_2$、$S_3$、$S_4$ 与 $S_5$，并将
        每个与邻近的 $2$ 的幂作比较，猜想
        $S_n$ 的公式。
  \item \emph{翻倍技巧}：写出和 $S_n$ 与
        $2 \times S_n$ 上下对齐，作差，并证明
        你的猜想：
        \[
        S_n = 1 + 2 + 4 + \dots + 2^{n-1} = 2^n - 1 .
        \]
  \item 国王一共许下了多少粒麦粒？用
        $2$ 的幂表示答案，并补全这个经典的
        论断：“整个棋盘上的麦粒比
        单独一个第六十五格所应放的少一粒。”
  \item 证明：棋盘后\omterm{def:g3:division:half}{一半}（第 $33$ 至
        $64$ 格）上的麦粒数\emph{恰好}为
        前\omterm{def:g3:division:half}{一半}的 $2^{32}$ 倍。
\end{enumerate}

\medskip
\noindent\textbf{第二部分——$2^{64}$ 有多大？}
$2^{10} = 1024 > 10^3$ 的比较（见
\cref{exo:g8:powers:11}）是下面所有估计的关键。
\begin{enumerate}[resume]
  \item 证明：
        $2^{64} = 2^4 \times \left(2^{10}\right)^6 > 1.6 \times
        10^{19}$.
  \item 一粒麦粒约重 $0.05$ g，即\
        $5 \times 10^{-2}$ g。证明所许诺的麦子
        重量大于 $8 \times 10^{17}$ g，并将其换算为
        吨（$1$ 吨 $= 10^6$ g）。
  \item 当今全世界每年约收获 $8 \times 10^8$
        吨小麦。国王许下的承诺，至少相当于
        全世界多少年的收成？
  \item 求使下式成立的最小正整数 $n$：
        $2^n > 10^6$ ——即翻倍多少次
        才能超过一百万。（精确计算 $2^{19}$ 与 $2^{20}$，
        利用 $2^{10} = 1024$。）
  \item 某骗子向你提议一个月的工资：第 $1$ 天给 $1$ 分，
        之后每天的工资为前一天的两倍。哪一天
        \emph{单日}工资本身首次超过一
        百万欧元（$10^8$ 分）？（精确计算 $2^{26}$ 与
        $2^{27}$。）
\end{enumerate}

\medskip
\noindent\textbf{第三部分——二进制砝码。}
一位商人有五枚砝码，分别重 $1$、$2$、$4$、$8$ 与 $16$ 克，
每种各一枚。她把其中一些放在天平的一只托盘上，
用以称量另一只托盘上的货物。
\begin{enumerate}[resume]
  \item 她要称出 $21$ g，应放哪些砝码？要称出
        $27$ g，又该放哪些？
  \item 说明为什么任何不小于 $16$ g 的目标\emph{必须}用到
        $16$ g 砝码，而任何不大于 $15$ g 的目标
        \emph{不得}用它。（第 3 题告诉你，
        砝码 $1, 2, 4, 8$ 至多能称出多少。）请说明同样的推理
        如何以次大的砝码重复，
        并在每个阶段都成立。
  \item 由此推出：从 $1$ 到 $31$ g 的每个整数目标都能称出，
        并且\emph{恰有一种称法}：每个介于 $1$ 与 $31$ 之间的数
        都是互不相同的 $2$ 的幂之和，
        且表示方式唯一。
  \item 商人又买了第六枚砝码，重 $32$ g。加上它后，她
        现在最多能称出多少克？把 $45$ g 写成
        互不相同的 $2$ 的幂之和。
  \item 在计算机中，一个“64 位”的数存放在 $64$
        个格子上，每格存放 $0$ 或 $1$ ——第 $k$
        格贡献 $2^{k-1}$，前提是它存放的是 $1$，与传说中的
        麦粒一样。利用第一部分，说明为什么
        这样一台机器能存储的整数恰好
        从 $0$ 到 $2^{64} - 1$。
\end{enumerate}
\end{problem}
